% Generated draft: requires mathematical and source review.
\documentclass[11pt]{article}
\usepackage[margin=25mm]{geometry}
\usepackage{fontspec}
\setmainfont{Latin Modern Roman}
\usepackage{amsmath,amssymb,mathtools,unicode-math}
\setmathfont{Latin Modern Math}
\usepackage{enumitem,needspace,xurl,hyperref}
\setlength{\parindent}{0pt}
\setlength{\parskip}{5pt}
\setlength{\emergencystretch}{3em}
\newcommand{\N}{\mathbb{N}}
\newcommand{\Z}{\mathbb{Z}}
\newcommand{\Q}{\mathbb{Q}}
\newcommand{\R}{\mathbb{R}}
\newcommand{\C}{\mathbb{C}}
\newcommand{\F}{\mathbb{F}}
\newcommand{\A}{\mathbb{A}}
\newcommand{\PP}{\mathbb{P}}
\newcommand{\E}{\mathbb{E}}
\newcommand{\eps}{\varepsilon}
\newcommand{\dd}{\,\mathrm{d}}
\begin{document}
\section*{UnsolvedMath ID: 30003259}
\textbf{Source problem code:} catalogue entry
\subsection{Definitions and mathematical statement}
Let \(R=\C[t]\), where \(t\) is an indeterminate, and let \(\mathcal L_{\mathrm{ring}}=\{0,1,+,-,\cdot\}\). For \(f\in R\setminus\{0\}\), let \(\deg(f)\) be its polynomial degree; in particular, every nonzero constant has degree \(0\). Define
\[
E^{\times}=\{(f,g)\in(R\setminus\{0\})^2:\deg(f)=\deg(g)\}.
\]
The source asks about a relation \(E\) on all of \(R\), but does not specify the degree of the zero polynomial. Thus its restriction to nonzero arguments is unambiguous, while its values on pairs involving zero require a convention. For example, adopting \(\deg(0)=-\infty\) would give \(E=E^{\times}\cup\{(0,0)\}\); this convention is not supplied by the record.

For a relation \(S\subseteq R^2\), definability in the ring language with parameters means that there exist a finite tuple \(\mathbf a\in R^r\) and an \(\mathcal L_{\mathrm{ring}}\)-formula \(\varphi(x,y,\mathbf u)\) such that
\[
\forall f,g\in R,\qquad (f,g)\in S\ \Longleftrightarrow\ R\models\varphi(f,g,\mathbf a).
\]
Under the usual convention allowing coefficients from \(R\), being Diophantine means that there exist integers \(m\geq0\), \(s\geq1\), and polynomials \(P_1,\ldots,P_s\in R[X,Y,Z_1,\ldots,Z_m]\) such that
\[
\forall f,g\in R,\qquad (f,g)\in S\ \Longleftrightarrow\ \exists\mathbf z\in R^m\ \bigwedge_{i=1}^{s}P_i(f,g,\mathbf z)=0.
\]
These finitely many coefficients constitute parameters. The source does not separately clarify whether its phrase ``with parameters'' also qualifies the Diophantine alternative.

\textbf{Degree relation question.} Is the degree-equivalence relation \(E\) on \(\C[t]\) Diophantine? Is it definable in \(\mathcal L_{\mathrm{ring}}\) with parameters? The intended relation agrees with \(E^{\times}\) on nonzero arguments; its treatment of zero, and the parameter convention for the Diophantine question, remain qualifications requiring review.

\textbf{Existential theory question.} Fix a prime \(p\), let \(\F_p\) be the field with \(p\) elements, and let \(K_p=\F_p(t)\) be its rational function field. Define
\[
\operatorname{Th}_{\exists}(K_p)=\left\{\exists x_1\cdots\exists x_n\,\psi(x_1,\ldots,x_n):\ n\geq0,\ \psi\text{ is quantifier-free in }\mathcal L_{\mathrm{ring}},\ K_p\models\exists x_1\cdots\exists x_n\,\psi\right\}.
\]
Is \(\operatorname{Th}_{\exists}(K_p)\) decidable? Equivalently, for this fixed \(p\), does an algorithm halt on every existential \(\mathcal L_{\mathrm{ring}}\)-sentence and correctly determine its truth in \(K_p\)? Here the language has no additional named parameters, in particular no constant symbol naming \(t\). The question is posed for each fixed prime; a single algorithm uniform in \(p\) is not explicitly requested.
\subsection{Short English statement}
The record asks whether equality of polynomial degrees over \(\C[t]\) is Diophantine or definable with parameters, and whether the parameter-free existential theory of \(\F_p(t)\) is decidable for a fixed prime \(p\). The treatment of the zero polynomial and the parameters allowed in the Diophantine alternative need clarification. The supplied open-status assessment is dated curation, not independent verification.
\subsection{Sources}
\begin{itemize}
\item Supplied local problem registry record: permanent UnsolvedMath ID \(30003259\), source problem code \texttt{catalogue entry}, title \emph{Degree Definability and Existential Theory of Rational Function Fields}; collection: Oberwolfach Reports Open Problems. No individual proposer is identified. The supplied proposed year is 2016.
\item Supplied source citation: \emph{Definability and Decidability Problems in Number Theory} (2017), Oberwolfach Reports, pp.\ 2793--2866, DOI \texttt{10.4171/owr/2016/49}. \url{https://doi.org/10.4171/owr/2016/49}. The report is described as documenting mathematical workshops and their open-problem discussions.
\item Supplied literature source, labelled ``Oberwolfach Reports 2017'': \url{https://ems.press/content/serial-article-files/46652}.
\item Literature assessment checked at \texttt{2026-08-21}: the supplied curation labels the problem open and states that the existential theory of \(\F_p(t)\) in the pure ring language remains open, and that decidability results with a named parameter do not resolve the parameter-free question. This is supplied dated curation derived from the OWR material, not independently verified here; its time-sensitive status and cited claims require independent verification. No separate literature assessment of the degree-equivalence question is supplied.
\item The registry labels the statement \texttt{corrected\_verified} and reports a manual semantic, typographic, and LaTeX review. These are supplied editorial assertions; this statement-only draft still requires human mathematical review.
\end{itemize}
\end{document}
